% Primitive: dynamic-programming — GUIDE-FIRST (plan 012 D-E): no parent
% series. This chapter partially delivers, in print, the wiki's standing
% promised row `ctc-alignment -> (dynamic programming as its own
% concept)`; the row stays promised until scenes exist. The video road's
% trellis scenes PERFORM the move without naming it, and are cited here
% only as users of the concept, never as teachers. Numbers trace to plan
% 001 anchors and small enumerable examples; problem answers are computed
% by answers/dynamic_programming.py.
\chapter{Dynamic programming}

\begin{narrative}
The last chapter ended on a sleight of hand worth catching red-handed. A
sum over more paths than anyone could list was priced exactly on a grid
of twenty circles, and the move was introduced in a single clause —
\emph{paths sharing a prefix share their count} — as if that were a
small bookkeeping remark. It is not small. It is one sighting of a
general method with its own name, its own conditions, and a habit of
turning impossible counts into short columns of additions. This chapter
takes the trellis back out and names what it did: \emph{dynamic
programming} — solve every small version of the problem once, store the
answers, and let the big version assemble itself from the store.

\section{The same question, asked over and over}

Start away from CTC, with the smallest example that misbehaves. The
Fibonacci numbers obey $F_n = F_{n-1} + F_{n-2}$ with $F_0 = 0$,
$F_1 = 1$. Computed exactly as written — each $F_n$ asking for its two
predecessors, each of those asking again — the questions repeat:
$F_{10}$ asks for $F_8$ twice (once itself, once through $F_9$), for
$F_7$ three times, for $F_6$ five times; the repetition compounds at
the same rate as the sequence itself. Answering $F_{10}$ this way takes
$177$ function calls to settle $11$ distinct questions.

The cure is almost insultingly simple: \emph{write the answers down}.
Compute $F_2, F_3, \dots$ in order, each from two stored numbers —
eleven entries, one addition each. The exponential call tree collapses
because two things were true at once. First, the subproblems
\emph{overlap}: the same question recurs, so a stored answer keeps
getting reused. Second, the stored answer is \emph{all the future
needs}: whoever asks for $F_8$ cares only about its value, never about
how the additions inside it were arranged. When both hold, the tree of
recomputation flattens into a table.

\section{The lattice, in its third costume}

The counting chapter closed on a walker crossing a lattice — $4$ rights
and $2$ ups, $\binom{6}{2} = 15$ routes, counted there by choosing
which steps go up. Count it again, the dynamic-programming way: label
every node with \emph{the number of routes that reach it}. A node can
only be entered from its left or lower neighbour, so its label is the
sum of two earlier labels — Pascal's addition, marching across the
grid. The corner lands on the same $15$, and no route was ever written
out. Each label is a stored subproblem: all routes agreeing on where
they currently stand are carried as one number from there on. The
exponential object (the list of routes) was reorganised into a small
table of counts — \emph{shared prefixes, counted once}.

\section{The trellis, named}

Now reread the last chapter with these words available. The mini
trellis from its third problem — transcript A over $T = 3$ frames,
states $(\varepsilon, A, \varepsilon)$ — enumerated $2^3 = 8$ raw paths
and tested each against the collapse map. The recurrence answered the
same question by filling nine cells:

\begin{center}
\begin{tikzpicture}[x=1.9cm, y=1.05cm]
  \foreach [count=\t from 1] \col in {{1,1,0},{1,2,1},{1,3,3}} {
    \foreach [count=\s from 1] \v in \col {
      \ifnum\v>0
        \node[circle, draw=Muted, minimum size=6.4mm, inner sep=0pt]
          (n\t\s) at (\t, -\s) {\small $\v$};
      \else
        \node[circle, draw=Muted!45, minimum size=6.4mm, inner sep=0pt]
          (n\t\s) at (\t, -\s) {};
      \fi
    }
    \node[above, color=CoolInk] at (\t, -0.55) {\small $t=\t$};
  }
  \foreach [count=\s from 1] \lab/\col in
    {$\varepsilon$/Muted, $A$/PaletteA, $\varepsilon$/Muted}
    \node[left, color=\col] at (0.55, -\s) {\lab};
  \foreach \t/\u in {1/2, 2/3} {
    \foreach \s/\r in {1/1, 2/2, 2/3, 3/3}
      \draw[Muted!60] (n\t\s) -- (n\u\r);
  }
  \draw[Muted!60] (n21) -- (n32);
  % The shared-prefix bundle: two partial paths end in this cell; their
  % count is stored once and every continuation reuses it.
  \draw[CoolInk, very thick] (n11) -- (n22);
  \draw[CoolInk, very thick] (n12) -- (n22);
  \node[circle, draw=AccentInk, very thick, minimum size=6.4mm,
    inner sep=0pt] at (2, -2) {};
  \node[below, color=CoolInk, align=center] at (2, -3.55)
    {\small two prefixes end here —\\[-2pt]\small stored once, reused after};
  \node[right, color=AccentInk] at (3.45, -2.5) {\small $3 + 3 = 6$};
\end{tikzpicture}
\end{center}

The highlighted cell is the whole method in one circle. Two partial
paths — $\varepsilon A$ and $AA$ — arrive at state $A$ after two
frames. From that cell onward their futures are identical: what a path
may legally do next depends only on \emph{where it stands}, not on how
it got there. So the two are stored as one number, and every
continuation is counted once instead of twice. That is the memo table
again: $\alpha_t(s)$ \emph{is} the stored answer to ``how much partial
weight reaches state $s$ by frame $t$?'', and the recurrence builds
each column from the stored one before it.

The same accounting, scaled: the AB trellis handled
$\anchor{001.raw81}$ raw paths, $\anchor{001.paths15}$ of them
accepted, on twenty cells. At $T = 100$, $U = 50$, the path count is
roughly $\anchor{001.astronomical}$ — and the trellis is
$100 \times 101$ cells, about ten thousand entries of three additions
and a multiplication each. The exponential object was never listed;
it was reorganised. Both of dynamic programming's conditions did real
work here: the subproblems overlap (astronomically many paths route
through each cell), and the cell is all the future needs (the collapse
map's legality — stay, advance, skip — reads only the current state).
Had the continuation rules depended on more of the past, each cell
would have had to remember more, and the grid would have grown states
until they did.

\section{Recognising a dynamic program in the wild}

The signature has two parts, and both must be present. There is a
brute-force count, sum, or best-choice over exponentially many
arrangements; and there is a small \emph{state} — the walker's node,
the trellis cell, the index $n$ — such that arrangements agreeing on
the state have interchangeable futures. Edit distance between two
strings (the state: a prefix of each), shortest routes through a road
network (the state: the junction you stand at), and the forward
algorithm this road met last chapter — the same reorganisation
runs hidden-state models generally~\cite{deep_learning-lawrence-r-rabiner-the-hmm-tutorial-proc-ieee-1989}
and is the standard route through
CTC~\cite{deep_learning-hannun-sequence-modeling-with-ctc-distill-2017}.
The boundary is just as useful to know: if subproblems never repeat,
a memo stores answers nobody asks for again — splitting a list into
disjoint halves is divide-and-conquer, not dynamic programming — and
if no small state summarises the past, the table itself grows
exponential and the method returns nothing.

\paragraph{Where this goes.} The trellis is not finished with this
road. When the gradient chapter differentiates the alignment sum, the
same grid will be swept \emph{backward} — a second dynamic program,
over suffixes — and the two sweeps together price every path through
every cell. Later, the logarithms chapter will find the trellis's
long products starving floating point, and the recurrence will move
into log-space to survive. And this chapter itself is a seed: the
video road performed dynamic programming without ever naming it, and
a future series owes the concept its own screen time — these pages
are its first draft.
\end{narrative}

\section*{Practice problems}

\begin{problem}
Computing $F_{10}$ by the bare recurrence $F_n = F_{n-1} + F_{n-2}$
(with $F_0 = 0$, $F_1 = 1$ answered directly), how many function calls
are made in total? How many \emph{distinct} questions were asked?
\begin{hint}
Count calls with their own recurrence: $c_n = 1 + c_{n-1} + c_{n-2}$,
with $c_0 = c_1 = 1$.
\end{hint}
\begin{solution}
The call count obeys $c_n = 1 + c_{n-1} + c_{n-2}$, which solves to
$c_n = 2F_{n+1} - 1$: the tree grows at the sequence's own rate. For
$n = 10$: $c_{10} = 2F_{11} - 1 = 2 \cdot 89 - 1 = 177$ calls — to
answer the $11$ distinct questions $F_0$ through $F_{10}$. The memo
table answers each distinct question once: eleven entries, one
addition each.
\end{solution}
\end{problem}

\begin{problem}
Label every node of the $4$-right, $2$-up lattice with the number of
routes reaching it, using only the addition rule. Confirm the far
corner agrees with the counting chapter's $\binom{6}{2}$.
\begin{hint}
Start with $1$ at the origin; every other node is the sum of its left
and lower neighbours (a missing neighbour contributes nothing).
\end{hint}
\begin{solution}
Filling row by row (bottom row first):
\[
\begin{matrix}
1 & 3 & 6 & 10 & 15\\
1 & 2 & 3 & 4 & 5\\
1 & 1 & 1 & 1 & 1
\end{matrix}
\]
The corner holds $15 = \binom{6}{2}$, matching the choose-the-ups
count — and matching $\anchor{001.paths15}$, the AB path count, a
third meeting with the same number: the trellis \emph{is} a lattice
walk with different legal steps.
\end{solution}
\end{problem}

\begin{problem}
Extend the mini trellis one more frame: for transcript A over $T = 4$
frames, compute the fourth column from $(1, 3, 3)$ and read off the
accepted count. Then confirm the count by a direct argument about
which raw paths collapse to A.
\begin{hint}
Stay and advance only — no skip exists on $(\varepsilon, A,
\varepsilon)$. For the direct argument: a path collapsing to A shows
one contiguous block of A's.
\end{hint}
\begin{solution}
Each cell adds its own state's previous count and the one above:
$(1,\ 3{+}1,\ 3{+}3) = (1, 4, 6)$, and the two accepting states give
$4 + 6 = 10$. Directly: a path collapsing to A is exactly one
contiguous block of A-frames inside $4$ frames — intervals of length
$1, 2, 3, 4$ number $4 + 3 + 2 + 1 = 10$. Two routes, one answer,
and the recurrence wrote three cells to the enumeration's sixteen
paths.
\end{solution}
\end{problem}

\begin{problem}[stretch]
At $T = 100$ frames and $U = 50$ letters, compare the work the two
routes perform: how many cells does the trellis fill, and how does
that stand against the raw repeat-free path count? What would happen
to the table if a path's legal continuations depended on its last
\emph{two} states instead of one?
\begin{hint}
The wrapped transcript has $2U + 1$ states. For the second part: what
must a cell now remember?
\end{hint}
\begin{solution}
The trellis fills $T (2U + 1) = 100 \times 101 = 10{,}100$ cells, a
handful of operations each, against roughly $\anchor{001.astronomical}$
paths — about $36$ orders of magnitude apart. If legality read the
last two states, ``where the path stands'' would no longer determine
its future: cells would need to be \emph{pairs} of states, squaring
the state count to order $10^4$ and the table to order $10^6$ —
still polynomial, but the lesson generalises: dynamic programming
pays exactly for how much past the future needs. When no small
summary exists, the table, not the path list, becomes the
exponential object.
\end{solution}
\end{problem}
